\documentclass[bild]{fiwthesis}  %% Bildunterschriften standardmäßig mit "Bild"

%% =====================
%%  Pakete (Empfehlungen)
%% =====================

\usepackage{lipsum}               
\usepackage{siunitx}              
\usepackage{textgreek}            
\usepackage{amsmath, amssymb}    
\usepackage{tabularray}           
\usepackage{graphicx}             
\usepackage{subcaption}           
\usepackage{pgfplots}             
\usepackage{listings}
\usepackage{xcolor}

\usepackage{tikz}
\usetikzlibrary{positioning, shapes, intersections}

\usepackage{setspace}
\usepackage[hidelinks]{hyperref}

\pgfplotsset{compat=1.18}

%% ===========
%%  Metadaten
%% ===========

\thesis{KI-System: Multimodaler Sprach-Assistent}
\title{Entwicklung eines multimodalen Sprach-Assistenten zur Verarbeitung gesprochener Sprache mittels Speech-to-Text, Natural Language Processing und visueller Informationsausgabe}
\date{18.04.2026}

\author{Ruba Shehadeh, Godfrey Uchechukwu Atulegwu, Kilian Ketelhohn}
\supervisor{Prof.~Dr.~Antje Raab-Düsterhöft}

\keywords{Python, Whisper, spaCy, Streamlit, NLP, Speech-to-Text, KI-Pipeline, Informationsextraktion, Bildgenerierung}

\makepdfmetadata  

%% ==========
%%  Präambel
%% ==========

\bibliography{quellen}              
%% =================================================================
%% verzeichnisse/glossar.tex
%% =================================================================

\newglossaryentry{crossplatform}{%
  name={Cross-Platform},
  description={Software-Entwicklungsansatz, bei dem Anwendungen mit einer gemeinsamen Codebasis auf mehreren Betriebssystemen (z. B. Android und iOS) ausführbar sind, hier umgesetzt mit Flutter.}
}

\newglossaryentry{speech2text}{%
  name={Speech-to-Text},
  description={Verfahren zur automatischen Umwandlung gesprochener Sprache in geschriebenen Text, im Projekt realisiert durch das Whisper-Modell.}
}

\newglossaryentry{nlp}{%
  name={Natural Language Processing},
  description={Teilgebiet der künstlichen Intelligenz zur Verarbeitung und Analyse natürlicher Sprache, z. B. für Tokenisierung, POS-Tagging und Informationsextraktion mit spaCy.}
}

\newglossaryentry{bootloader}{%
  name={Bootstrapping / Pipeline-Start},
  description={Initialisierung der Verarbeitungskette im System, bei der Audioaufnahme, Transkription und Analyse sequenziell gestartet werden.}
}

\newglossaryentry{relational}{%
  name={Relationales Modell},
  description={Datenmodell, bei dem Informationen strukturiert in Tabellen gespeichert und über Primär- und Fremdschlüssel (z. B. in SQL-Datenbanken) miteinander verknüpft werden.}
}

\newglossaryentry{offlinefirst}{%
  name={Offline-First},
  description={Architekturansatz, bei dem die Anwendung vollständig lokal funktionsfähig ist und Daten bei bestehender Verbindung synchronisiert werden.}
}

\newglossaryentry{widget}{%
  name={Widget},
  description={Zentrale UI-Baueinheit in Flutter. Jede visuelle oder interaktive Komponente der Streamlit- bzw. Flutter-Oberfläche basiert auf Widgets.}
}

\newglossaryentry{pipeline}{%
  name={KI-Pipeline},
  description={Modulare Verarbeitungskette aus Speech-to-Text, NLP-Analyse und visueller Ausgabe zur schrittweisen Transformation von Sprache in strukturierte Informationen.}
}

% =========================================================
% Dual Entries (Akronyme + Definitionen)
% =========================================================

\newdualentry{api}
{API}
{Application Programming Interface}
{Programmierschnittstelle zur standardisierten Kommunikation zwischen Systemkomponenten, insbesondere für die Anbindung externer Dienste und Bilddatenquellen (z. B. Unsplash API).}

\newdualentry{json}
{JSON}
{JavaScript Object Notation}
{Textbasiertes Datenformat zur strukturierten Übertragung von Informationen zwischen Client (Streamlit-App) und externen oder lokalen Services.}

\newdualentry{stt}
{STT}
{Speech-to-Text}
{Automatische Umwandlung von Audioeingaben in Textdaten mittels Whisper-Modell als erster Schritt der KI-Pipeline.}       
%% =================================================================
%% verzeichnisse/abkuerzungen.tex
%% =================================================================

% Nutzungshinweise:
% \gls{tag}     -> Erstes Auftreten: Vollform (Kurzform)
% \glspl{tag}   -> Pluralform
% \acrfull{tag} -> Erzwingt Vollform

% =========================================================
% KI / Sprachverarbeitung / Frameworks
% =========================================================
\newacronym{stt}{STT}{Speech-to-Text}
\newacronym{nlp}{NLP}{Natural Language Processing}
\newacronym{llm}{LLM}{Large Language Model}
\newacronym{api}{API}{Application Programming Interface}
\newacronym{rest}{REST}{Representational State Transfer}
\newacronym{json}{JSON}{JavaScript Object Notation}
\newacronym{sdk}{SDK}{Software Development Kit}

% =========================================================
% Datenverarbeitung & Backend
% =========================================================
\newacronym{sql}{SQL}{Structured Query Language}
\newacronym{sqlite}{SQLite}{Serverlose, eingebettete SQL-Datenbank-Engine}
\newacronym{etl}{ETL}{Extract, Transform, Load}
\newacronym{pipeline}{Pipeline}{Modulare Verarbeitungskette zur schrittweisen Datenverarbeitung}

% =========================================================
% KI-Modelle & Verarbeitung
% =========================================================
\newacronym{whisper}{Whisper}{Automatisches Speech-to-Text-Modell von OpenAI}
\newacronym{spacy}{spaCy}{Bibliothek zur industriellen Verarbeitung natürlicher Sprache}
\newacronym{token}{Token}{Einzelnes linguistisches Element innerhalb eines Textes}

% =========================================================
% UI / UX / Frontend (Flutter / Streamlit)
% =========================================================
\newacronym{ui}{UI}{User Interface}
\newacronym{ux}{UX}{User Experience}
\newacronym{gui}{GUI}{Graphical User Interface}
\newacronym{flutter}{Flutter}{Cross-Platform UI Framework von Google}
\newacronym{widget}{Widget}{Grundbaustein der Benutzeroberfläche in Flutter}

% =========================================================
% Mobile / System
% =========================================================
\newacronym{apk}{APK}{Android Application Package}
\newacronym{os}{OS}{Operating System}
\newacronym{ram}{RAM}{Random Access Memory}
\newacronym{cpu}{CPU}{Central Processing Unit}
\newacronym{gpu}{GPU}{Graphics Processing Unit}
\newacronym{wlan}{WLAN}{Wireless Local Area Network}

% =========================================================
% Audio / Datenquellen
% =========================================================
\newacronym{mic}{MIC}{Mikrofon als Audioeingabegerät}
\newacronym{audio}{Audio}{Digitale Schallsignaldaten als Eingabequelle für STT-Systeme}  
%% =================================================================
%% verzeichnisse/symbole.tex
%% =================================================================

\newglossaryentry{symb:alpha}{
  name={$\alpha$},
  description={Rotationswinkel der 3D-Visualisierung um die vertikale Achse (z. B. Kamera-Azimuth im Modellviewer der Anwendung).},
  sort=symbolalpha
}

\newglossaryentry{symb:beta}{
  name={$\beta$},
  description={Neigungswinkel der 3D-Kamera bzw. des dargestellten Objekts im interaktiven Visualisierungssystem.},
  sort=symbolbeta
}

\newglossaryentry{symb:delta_t}{
  name={$\Delta t$},
  description={Zeitintervall zwischen zwei aufeinanderfolgenden Verarbeitungsschritten innerhalb der KI-Pipeline (z. B. STT → NLP → Visualisierung).},
  sort=symboldeltat
}

\newglossaryentry{symb:n}{
  name={$n$},
  description={Anzahl der verarbeiteten oder lokal gespeicherten Datensätze (z. B. transkribierte Audiosegmente oder erkannte Entitäten).},
  sort=symboln
}

\newglossaryentry{symb:fps}{
  name={$\mathrm{FPS}$},
  description={Bildwiederholrate der grafischen Ausgabe (z. B. Streamlit-Interface oder 3D-Renderer in Flutter) in Frames per Second.},
  sort=symbolfps
}

\newglossaryentry{symb:progress}{
  name={$P$},
  description={Fortschrittsvariable der Verarbeitungspipeline, normiert im Bereich $0 \leq P \leq 1$, z. B. für Transkriptions- oder Analysefortschritt.},
  sort=symbolp
}

\newglossaryentry{symb:confidence}{
  name={$C$},
  description={Konfidenzwert eines KI-Modells (z. B. Whisper oder spaCy) für eine erkannte Ausgabe oder Klassifikation.},
  sort=symbolc
}

\newglossaryentry{symb:tokens}{
  name={$T$},
  description={Anzahl der linguistischen Tokens innerhalb eines Textes, die in der NLP-Verarbeitung analysiert werden.},
  sort=symbolt
}        

\makeglossaries{}                   
\makeindex[
  intoc=true,
  title=Index,
  columns=2
]{}
\indexsetup{headers={\indexname}{\indexname}}

%% ===============
%%  Beginn Thesis
%% ===============

\begin{document}

%% ============
%%  Vorspann
%% ============

\maketitle

% Aufgabenstellung
\maketask{
Die Aufgabe dieser Arbeit besteht in der Entwicklung eines multimodalen KI-Systems, das gesprochene Sprache verarbeitet, linguistisch analysiert und die extrahierten Informationen visuell darstellt.

Das System nimmt Audioeingaben über ein Mikrofon auf, wandelt diese mithilfe von Speech-to-Text (Whisper) in Text um und analysiert diesen anschließend mit Natural Language Processing (spaCy).

Besonderer Fokus liegt auf der Identifikation von Verben und Nomen, die als Grundlage für eine visuelle Repräsentation dienen. Diese werden genutzt, um passende Bilder über eine Bilddatenbank (Unsplash API) abzurufen und im Benutzerinterface darzustellen.

Ziel ist die Entwicklung einer modularen KI-Pipeline, die Sprache strukturiert verarbeitet und visuell interpretiert.
}

\maketoc[compact]  

%% ==========
%%  Textteil
%% ==========

%% kapitel/einleitung.tex
\chapter{Einleitung}
\label{chap:einleitung}

\subsection{Motivation}
Die Interaktion zwischen Mensch und Maschine hat sich in den letzten Jahren grundlegend gewandelt. Klassische Eingabemethoden wie Tastatur und Maus werden zunehmend durch natürlichere Kommunikationsformen ergänzt. Insbesondere die Verarbeitung gesprochener Sprache ermöglicht eine intuitive und barrierearme Nutzung digitaler Systeme.

Moderne Anwendungen erfordern jedoch nicht nur die reine Umwandlung von Sprache in Text, sondern die Fähigkeit, Informationen kontextuell zu verstehen, zu strukturieren und visuell aufzubereiten. Genau hier setzt das Konzept eines \textit{multimodalen Sprachassistenten} an: Durch die Kombination verschiedener Modalitäten – wie Sprache, Text und visuelle Darstellung – entsteht ein interaktives System, das komplexe Inhalte verständlich und nutzerfreundlich vermittelt.

Die Herausforderung besteht darin, diese unterschiedlichen Technologien – darunter Speech-to-Text, Natural Language Processing und visuelle Informationsdarstellung – in einem integrierten System zu vereinen. Ziel ist es, eine Anwendung zu entwickeln, die nicht nur Sprache erkennt, sondern daraus aktiv verwertbare Informationen generiert und diese anschaulich präsentiert.

\subsection{Zielsetzung des Projekts}
Das Ziel dieser Arbeit ist die Konzeption und Entwicklung eines intelligenten, multimodalen Sprachassistenten auf Basis moderner Softwaretechnologien. Das System soll als interaktive Schnittstelle zwischen Nutzer und digitalen Inhalten fungieren und dabei verschiedene Eingabe- und Ausgabekanäle kombinieren.

Im Mittelpunkt steht die Entwicklung eines prototypischen Systems, das gesprochene Sprache in Echtzeit verarbeitet, analysiert und in strukturierte sowie visuelle Informationen überführt. Dabei werden verschiedene Komponenten zu einer durchgängigen Verarbeitungskette verbunden.

Zu den zentralen Zielen gehören:
\begin{itemize}
\item \textbf{Sprachverarbeitung:} Implementierung einer Speech-to-Text-Pipeline zur Umwandlung gesprochener Sprache in Textdaten.
\item \textbf{Sprachverständnis:} Einsatz von Natural Language Processing zur Extraktion relevanter Informationen und semantischer Strukturen.
\item \textbf{Datenverarbeitung:} Transformation der extrahierten Inhalte in strukturierte Datenformate zur Weiterverarbeitung.
\item \textbf{Visualisierung:} Darstellung der gewonnenen Informationen in Form von interaktiven und grafischen Elementen.
\item \textbf{Systemintegration:} Zusammenführung aller Komponenten in einer webbasierten Anwendung mit intuitiver Benutzeroberfläche.
\end{itemize}

\subsection{Struktur der Dokumentation}
Die vorliegende Arbeit beschreibt den vollständigen Entwicklungsprozess des multimodalen Sprachassistenten. Beginnend mit den theoretischen Grundlagen der Sprachverarbeitung und Datenanalyse werden anschließend die methodischen Ansätze sowie die Systemarchitektur erläutert.

Darauf aufbauend folgt die detaillierte Beschreibung der Implementierung, einschließlich der eingesetzten Technologien und zentralen Funktionen des Systems. Abschließend werden die Ergebnisse evaluiert und mögliche Erweiterungen sowie zukünftige Entwicklungsperspektiven diskutiert.
          
\chapter{Grundlagen}
\label{sec:grundlagen}

\section{Multimodale KI-Systeme}
Multimodale KI-Systeme kombinieren verschiedene Arten der Datenverarbeitung, beispielsweise Sprache, Text und visuelle Informationen. Ziel solcher Systeme ist es, Informationen aus einer Eingabemodalität (z.B. gesprochene Sprache) in eine andere Modalität (z.B. strukturierter Text oder Bild) zu überführen.

Im Rahmen dieser Arbeit wird eine KI-Pipeline entwickelt, die gesprochene Sprache verarbeitet, linguistisch analysiert und anschließend visuell repräsentiert.

\section{Speech-to-Text mit Whisper}
Die Grundlage des Systems bildet die automatische Spracherkennung (Speech-to-Text).

\textbf{OpenAI Whisper} ist ein neuronales Netzwerkmodell zur robusten Transkription gesprochener Sprache in Text. Es basiert auf einem Encoder-Decoder-Transformer-Ansatz und wurde auf großen Mengen multilingualer Audiodaten trainiert.

Wesentliche Eigenschaften:
\begin{itemize}
    \item Hohe Genauigkeit bei verschiedenen Akzenten und Sprachstilen
    \item Unterstützung mehrerer Sprachen, darunter Deutsch
    \item Fähigkeit zur Verarbeitung von verrauschten Audioaufnahmen
\end{itemize}

\section{Natural Language Processing (NLP)}
Natural Language Processing (NLP) bezeichnet die Verarbeitung und Analyse natürlicher Sprache durch Computer.

Im Projekt wird hierfür das Framework \textbf{spaCy} verwendet.

\subsection{Tokenisierung}
Die Tokenisierung beschreibt die Zerlegung eines Satzes in einzelne Einheiten (Tokens), meist Wörter oder Satzzeichen. Diese Strukturierung ist die Grundlage für alle weiteren Analyseprozesse.

\subsection{Part-of-Speech Tagging}
Beim Part-of-Speech (POS) Tagging werden jedem Token grammatikalische Kategorien zugewiesen, beispielsweise:
\begin{itemize}
    \item \textbf{VERB} (Verben / Handlungen)
    \item \textbf{NOUN} (Nomen / Objekte)
\end{itemize}

Diese Informationen sind entscheidend für die spätere Informationsextraktion.

\section{Informationsextraktion}
Die Informationsextraktion beschreibt den Prozess, relevante Inhalte aus unstrukturiertem Text zu gewinnen.

Im Kontext dieses Projekts werden insbesondere:
\begin{itemize}
    \item Verben zur Analyse von Handlungen
    \item Nomen zur Beschreibung von Objekten und Entitäten
\end{itemize}

extrahiert und weiterverarbeitet.

\section{Semantische Bildzuordnung}
Ein zentrales Konzept des Systems ist die semantische Verknüpfung von Sprache und visuellen Inhalten.

Dabei werden extrahierte Nomen als Suchbegriffe verwendet, um passende Bilder aus einer Bilddatenbank (z.B. Unsplash API) abzurufen.

Dieser Ansatz stellt eine einfache Form der semantischen Abbildung dar, bei der sprachliche Konzepte in visuelle Repräsentationen überführt werden.

\section{Streamlit als Benutzeroberfläche}
\textbf{Streamlit} ist ein Python-Framework zur schnellen Entwicklung interaktiver Webanwendungen für Daten- und KI-Projekte.

Im Projekt übernimmt Streamlit folgende Aufgaben:
\begin{itemize}
    \item Darstellung der Texteingabe und Audioeingabe
    \item Visualisierung der NLP-Ergebnisse
    \item Anzeige generierter oder abgerufener Bilder
\end{itemize}

\section{Architektur von KI-Pipelines}
Die Anwendung folgt dem Prinzip einer modularen KI-Pipeline.

Diese besteht aus klar getrennten Verarbeitungsschritten:
\begin{enumerate}
    \item Audioeingabe (Mikrofon)
    \item Speech-to-Text (Whisper)
    \item Textanalyse (spaCy)
    \item Informationsextraktion
    \item Visuelle Ausgabe (Bildsuche)
\end{enumerate}

Diese modulare Struktur ermöglicht eine einfache Erweiterbarkeit und Austauschbarkeit einzelner Komponenten.

\section{Zusammenfassung für Fachfremde}
Das System kann vereinfacht als ein intelligenter Sprach-zu-Bild-Übersetzer verstanden werden. Gesprochene Sprache wird zunächst in Text umgewandelt, anschließend analysiert und schließlich in passende Bilder übersetzt.

Dadurch wird Sprache nicht nur lesbar, sondern auch visuell erfahrbar gemacht.          
\chapter{Methodik}
\label{sec:methodik}

\section{Vorgehensweise}
Die Entwicklung des \textit{multimodalen Sprach-Assistenten} folgt einem modularen, ingenieurwissenschaftlich orientierten Ansatz zur Verarbeitung natürlicher Sprache. Ziel ist die Transformation unstrukturierter Sprachdaten in strukturierte und visuell interpretierbare Informationen.

Die Vorgehensweise gliedert sich in folgende Phasen:

\begin{enumerate}
    \item \textbf{Anforderungsanalyse und Use-Case-Definition:}  
    Analyse typischer Anwendungsszenarien von Sprach-KI-Systemen. Im Fokus steht die Frage, wie gesprochene Sprache in Echtzeit verarbeitet, analysiert und visuell dargestellt werden kann.

    \item \textbf{Pipeline-Architektur (Modulares KI-System):}  
    Entwurf einer klar strukturierten Verarbeitungskette bestehend aus Speech-to-Text, NLP-Analyse, Informationsextraktion und visueller Ausgabe. Die Architektur folgt dem Prinzip der modularen Austauschbarkeit einzelner Komponenten.

    \item \textbf{Prototyping und UI-Entwicklung:}  
    Implementierung einer interaktiven Weboberfläche mit Streamlit, welche Audio- und Texteingaben entgegennimmt und die Ergebnisse der KI-Pipeline in Echtzeit visualisiert.
\end{enumerate}

\section{Konzept der KI-Pipeline}
Im Gegensatz zu monolithischen KI-Systemen basiert die Umsetzung auf einer \textbf{sequenziellen Verarbeitungspipeline}, bei der jeder Schritt klar definierte Aufgaben übernimmt.

\subsection{Pipeline-Struktur}
Die Architektur besteht aus folgenden aufeinander aufbauenden Komponenten:

\begin{enumerate}
    \item \textbf{Audio Input Layer:} Aufnahme oder Upload von Spracheingaben.
    \item \textbf{Speech-to-Text Layer:} Transkription mittels OpenAI Whisper.
    \item \textbf{NLP Layer:} Linguistische Analyse mit spaCy.
    \item \textbf{Information Extraction Layer:} Extraktion von Verben und Nomen.
    \item \textbf{Visualization Layer:} Semantische Abbildung auf Bilder.
\end{enumerate}

\subsection{Designentscheidungen der Pipeline}
\begin{itemize}
    \item \textbf{Modularität:} Jede Stufe der Pipeline ist unabhängig implementiert und kann einzeln ersetzt oder erweitert werden.
    \item \textbf{Echtzeitfähigkeit:} Die Architektur ist darauf ausgelegt, Eingaben ohne komplexe Zwischenspeicherung zu verarbeiten.
    \item \textbf{Erweiterbarkeit:} Neue Modelle (z.B. Embedding-basierte oder generative KI-Modelle) können leicht integriert werden.
\end{itemize}

\section{Datenverarbeitung und semantisches Mapping}
Ein zentrales methodisches Konzept ist die Transformation von Sprache in semantische Schlüsselbegriffe.

\subsection{Tokenisierung und linguistische Analyse}
Der transkribierte Text wird mithilfe von spaCy in einzelne Tokens zerlegt und grammatikalisch annotiert (Part-of-Speech Tagging). Dadurch können strukturelle Elemente wie Verben und Nomen systematisch identifiziert werden.

\subsection{Semantische Reduktion}
Die extrahierten Nomen werden auf ihre semantische Relevanz reduziert, um als Suchbegriffe für die Bildgenerierung oder Bildsuche zu dienen. Dieser Prozess stellt eine vereinfachte Form des Semantic Mapping dar.

\section{Systemarchitektur}
Methodisch basiert das System auf einer klar getrennten Schichtenarchitektur:

\begin{enumerate}
    \item \textbf{Presentation Layer:} Streamlit Web-Interface
    \item \textbf{Processing Layer:} Whisper + spaCy Pipeline
    \item \textbf{Application Logic Layer:} Extraktion und Mapping-Logik
    \item \textbf{Output Layer:} Textvisualisierung und Bilddarstellung
\end{enumerate}

\section{Zusammenfassung}
Die gewählte Methodik kombiniert Konzepte aus der Sprachverarbeitung, der maschinellen Linguistik und der modularen Softwarearchitektur. Durch die konsequente Trennung der Verarbeitungsschritte entsteht ein flexibles und erweiterbares KI-System, das Sprache systematisch in strukturierte und visuelle Informationen überführt.            
\chapter{Praktische Umsetzung: Implementierung}
\label{sec:implementierung}

\section{Einleitung}
In diesem Kapitel wird die technische Umsetzung des \textit{multimodalen Sprach-Assistenten} beschrieben. Der Fokus liegt auf der Integration der einzelnen Komponenten der KI-Pipeline, bestehend aus Speech-to-Text, Natural Language Processing sowie der visuellen Ausgabe der analysierten Inhalte.

Ziel der Implementierung ist die Realisierung eines durchgängigen Systems, das gesprochene Sprache verarbeitet, strukturiert analysiert und in visuelle Informationen überführt.

\section{Workflow der Implementierung}
Die Entwicklung erfolgte modular entlang der einzelnen Stufen der KI-Pipeline:

\begin{enumerate}
    \item \textbf{Audioeingabe:} Aufnahme oder Upload von Sprachdaten über die Benutzeroberfläche (Streamlit).
    
    \item \textbf{Speech-to-Text Verarbeitung:} Umwandlung der Audiodaten in Text mittels \textbf{OpenAI Whisper}.
    
    \item \textbf{NLP-Verarbeitung:} Analyse des transkribierten Textes mit \textbf{spaCy}, inklusive Tokenisierung und Part-of-Speech Tagging.
    
    \item \textbf{Informationsextraktion:} Identifikation relevanter sprachlicher Elemente, insbesondere Verben und Nomen.
    
    \item \textbf{Visuelle Ausgabe:} Nutzung extrahierter Nomen als Suchbegriffe für eine Bilddatenbank (z.\,B. Unsplash API).
    
    \item \textbf{UI-Darstellung:} Visualisierung der Ergebnisse innerhalb einer Streamlit-Webanwendung.
\end{enumerate}

\section{Integration der Systemkomponenten}

\subsection{Speech-to-Text Modul}
Das Speech-to-Text Modul basiert auf dem Whisper-Modell, welches Audiosignale in Textsequenzen umwandelt. Die Implementierung erfolgt so, dass sowohl Audio-Dateien als auch potenziell Live-Eingaben verarbeitet werden können.

\subsection{NLP-Analyse Pipeline}
Die Verarbeitung des Textes erfolgt über spaCy. Dabei wird der Text zunächst tokenisiert und anschließend grammatikalisch analysiert.

Die wichtigsten Ergebnisse sind:
\begin{itemize}
    \item Identifikation von Verben (Handlungen)
    \item Extraktion von Nomen (Objekte)
    \item Strukturierung des Satzes in linguistische Einheiten
\end{itemize}

\subsection{Informationsextraktion und Mapping}
Die extrahierten Nomen werden als semantische Schlüsselbegriffe genutzt. Diese werden in Suchanfragen transformiert, um passende visuelle Inhalte zu finden.

Dieser Schritt stellt eine einfache Form der semantischen Abbildung zwischen Sprache und Bild dar.

\subsection{Benutzeroberfläche (Streamlit)}
Die gesamte Interaktion erfolgt über eine Streamlit-Webanwendung. Diese übernimmt:

\begin{itemize}
    \item Texteingabe und Audio-Upload
    \item Darstellung der NLP-Ergebnisse
    \item Visualisierung der markierten Verben im Text
    \item Anzeige der generierten bzw. abgerufenen Bilder
\end{itemize}

\section{Praktische Umsetzungsschritte}
Die Implementierung wurde in einzelne Python-Module unterteilt:

\begin{itemize}
    \item \texttt{speech\_to\_text.py}: Verarbeitung von Audioeingaben mit Whisper
    \item \texttt{nlp\_analysis.py}: Textanalyse mit spaCy (Tokenisierung, POS-Tagging)
    \item \texttt{image\_generator.py}: Generierung von Bild-URLs basierend auf Nomen
    \item \texttt{app.py}: Streamlit-Frontend und Orchestrierung der Pipeline
\end{itemize}

\section{Besondere Herausforderungen und Lösungen}

\begin{itemize}
    \item \textbf{Modellabhängigkeiten:} Die Integration von Whisper erforderte eine stabile Python-Umgebung mit kompatiblen Versionen von NumPy und PyTorch.
    
    \item \textbf{Audioverarbeitung:} Unterschiedliche Audioformate wurden durch temporäre Konvertierung in ein einheitliches Format gelöst.
    
    \item \textbf{Semantische Genauigkeit:} Die Qualität der Bildzuordnung hängt stark von der Auswahl relevanter Nomen ab, weshalb eine Filterung der Token implementiert wurde.
\end{itemize}

\section{Testing und Qualitätssicherung}
Die Qualitätssicherung erfolgte durch manuelle und funktionale Tests der Pipeline:

\begin{itemize}
    \item Test verschiedener Sprachinputs (einfach, komplex, mehrsatzig)
    \item Validierung der korrekten POS-Erkennung durch spaCy
    \item Überprüfung der Bildzuordnung anhand semantischer Konsistenz
    \item UI-Tests innerhalb der Streamlit-Anwendung
\end{itemize}

\section{Zusammenfassung}
Die Implementierung zeigt die erfolgreiche Realisierung einer modularen KI-Pipeline. Durch die klare Trennung der einzelnen Verarbeitungsschritte konnte ein flexibles und erweiterbares System geschaffen werden, das Spracheingaben in strukturierte und visuelle Informationen überführt.     
\chapter{Evaluation}
\label{sec:evaluation}

\section{Testumgebung und Datensätze}
Die Evaluation des \textit{multimodalen Sprach-Assistenten} erfolgte auf einem lokalen Entwicklungsrechner unter macOS sowie innerhalb einer isolierten Python-Umgebung (venv).

Zur funktionalen Überprüfung des Systems wurden verschiedene Sprach- und Texteingaben verwendet, die unterschiedliche linguistische Strukturen abbilden. Diese umfassen einfache Sätze, komplexe Mehrsatzstrukturen sowie handlungsorientierte Beschreibungen, um die Robustheit der NLP-Pipeline zu testen.

\section{Performance-Analyse}

\subsection{Speech-to-Text (Whisper)}
Die Spracherkennung mittels OpenAI Whisper zeigte eine hohe Genauigkeit bei klar gesprochenen Eingaben.

\begin{itemize}
    \item Gute Erkennungsrate bei deutschen Standardsätzen
    \item Leichte Genauigkeitsverluste bei stark verrauschten Audioeingaben
    \item Transkriptionszeit abhängig von Modellgröße (base vs. small)
\end{itemize}

Die Verarbeitung erfolgt nahezu in Echtzeit für kurze Audiosequenzen, wodurch eine interaktive Nutzung ermöglicht wird.

\subsection{Natural Language Processing (spaCy)}
Die linguistische Analyse durch spaCy liefert stabile Ergebnisse im Bereich Tokenisierung und Part-of-Speech Tagging.

\begin{itemize}
    \item Verben werden zuverlässig als Handlungen erkannt
    \item Nomen werden konsistent als Objekte extrahiert
    \item Satzstruktur bleibt auch bei längeren Eingaben stabil analysierbar
\end{itemize}

Die Performance der NLP-Komponente ist unabhängig von der Eingabelänge für typische Nutzungsszenarien ausreichend.

\subsection{Bildgenerierung und semantische Abbildung}
Die semantische Umwandlung von Sprache in visuelle Inhalte erfolgt über die Extraktion von Nomen, welche als Suchbegriffe für eine Bilddatenquelle verwendet werden.

\begin{itemize}
    \item Hohe Relevanz bei einfachen Objektbegriffen (z.B. "Hund", "Park")
    \item Eingeschränkte Genauigkeit bei abstrakten Begriffen
    \item Schnelle Bildauslieferung über externe API (Unsplash)
\end{itemize}

\section{Funktionale Ergebnisse}

Die funktionale Evaluation des Gesamtsystems ergab folgende Erkenntnisse:

\begin{itemize}
    \item \textbf{End-to-End Pipeline:} Die vollständige Verarbeitungskette (Audio $\rightarrow$ Text $\rightarrow$ NLP $\rightarrow$ Bild) funktioniert stabil und ohne manuelle Eingriffe.
    
    \item \textbf{Textvisualisierung:} Die Hervorhebung von Verben im Text verbessert die Nachvollziehbarkeit der linguistischen Analyse deutlich.

    \item \textbf{Interaktive Nutzung:} Die Streamlit-Oberfläche ermöglicht eine einfache Bedienung und unmittelbare Rückmeldung der KI-Ergebnisse.

    \item \textbf{Robustheit:} Das System bleibt auch bei fehlerhaften oder unvollständigen Eingaben stabil und liefert interpretierbare Ergebnisse.
\end{itemize}

\section{Zusammenfassende Beurteilung}
Die Evaluation zeigt, dass der entwickelte multimodale Sprach-Assistent die definierten Anforderungen erfüllt. Besonders hervorzuheben ist die erfolgreiche Kombination aus Spracherkennung, linguistischer Analyse und visueller Repräsentation.

Die modulare Architektur erlaubt eine einfache Erweiterung des Systems, beispielsweise durch den Einsatz leistungsfähigerer Sprachmodelle oder generativer Bild-KI. Insgesamt stellt das Projekt eine funktionale KI-Pipeline dar, die Sprache in strukturierte und visuell interpretierbare Informationen überführt.          
\chapter{Fazit und Ausblick}
\label{sec:fazit}

Die vorliegende Arbeit zeigt, wie moderne KI-Technologien zur Verarbeitung gesprochener Sprache in einer durchgängigen multimodalen Pipeline kombiniert werden können. Der entwickelte \textit{multimodale Sprach-Assistent} transformiert Spracheingaben in strukturierte Textinformationen und erweitert diese um eine visuelle Repräsentation.

\section{Zusammenfassung der Ergebnisse}
Im Rahmen des Projekts wurden folgende zentrale Ziele erreicht:

\begin{itemize}
    \item \textbf{Speech-to-Text Verarbeitung:} Durch den Einsatz von OpenAI Whisper wurde eine zuverlässige Umwandlung gesprochener Sprache in Text realisiert, die als Grundlage der weiteren Verarbeitung dient.

    \item \textbf{Linguistische Analyse:} Mithilfe von spaCy konnte der transkribierte Text erfolgreich tokenisiert und grammatikalisch analysiert werden. Insbesondere die Identifikation von Verben und Nomen ermöglicht eine strukturierte Interpretation der Inhalte.

    \item \textbf{Informationsextraktion und Visualisierung:} Die extrahierten Nomen werden zur semantischen Bildsuche genutzt und ermöglichen die visuelle Darstellung sprachlicher Inhalte über externe Bildquellen.

    \item \textbf{Benutzeroberfläche:} Die Umsetzung mit Streamlit ermöglicht eine interaktive Nutzung des Systems, bei der Spracheingaben, Analyseergebnisse und Bilder in Echtzeit dargestellt werden.

    \item \textbf{Pipeline-Architektur:} Die modulare Struktur der Anwendung erlaubt eine klare Trennung der Verarbeitungsschritte von Audioeingabe bis zur visuellen Ausgabe.
\end{itemize}

\section{Persönliches Resümee}
Die Entwicklung des Systems ermöglichte tiefere Einblicke in die praktische Anwendung moderner KI-Technologien. Besonders die Kombination aus Spracherkennung, NLP und visueller Repräsentation verdeutlicht die Komplexität multimodaler Systeme.

Herausfordernd war insbesondere die Abstimmung der einzelnen Module innerhalb einer stabilen End-to-End-Pipeline sowie die effiziente Verarbeitung der Daten in einer interaktiven Webanwendung.

\section{Ausblick}
Trotz des erreichten Funktionsumfangs bietet das System zahlreiche Erweiterungsmöglichkeiten:

\begin{itemize}
    \item \textbf{Integration generativer KI-Modelle:} Anstelle einfacher Bildsuche könnte ein Bildgenerierungsmodell (z.\,B. DALL·E) eingesetzt werden, um vollständig neue visuelle Inhalte zu erzeugen.

    \item \textbf{Erweiterte semantische Analyse:} Der Einsatz von Embedding-Modellen könnte die semantische Interpretation von Texten deutlich verbessern und abstrakte Begriffe besser abbilden.

    \item \textbf{Echtzeit-Sprachverarbeitung:} Eine Live-Mikrofonverarbeitung würde eine noch natürlichere Interaktion mit dem System ermöglichen.

    \item \textbf{Erweiterte KI-Interaktion:} Die Integration eines dialogbasierten LLM könnte das System zu einem vollständigen intelligenten Assistenten erweitern.
\end{itemize}

Abschließend lässt sich festhalten, dass der entwickelte multimodale Sprach-Assistent eine funktionale KI-Pipeline darstellt, die Sprache erfolgreich in strukturierte und visuell interpretierbare Informationen überführt und damit eine solide Grundlage für zukünftige multimodale KI-Anwendungen bildet.               

%% =========
%%  Anlagen
%% =========

\begin{appendices}
  \chapter{Quellcode-Dokumentation}
  %% ===============================
%% anlagen/listings.tex
%% Quellcode-Auszüge: Multimodaler Sprach-Assistent
%% ===============================

% WICHTIG:
% KEIN \chapter hier! Wird in der Hauptdatei definiert.

% ===============================
% Listings Setup für Python
% ===============================

\lstdefinelanguage{Python}{
    morekeywords={as,assert,break,class,continue,def,del,elif,else,except,False,finally,for,from,global,if,import,in,is,lambda,None,nonlocal,not,or,pass,raise,return,True,try,while,with,yield},
    sensitive=true,
    morecomment=[l]{\#},
    morestring=[b]',
    morestring=[b]"
}

\lstset{
    language=Python,
    basicstyle=\ttfamily\small,
    numbers=left,
    numberstyle=\tiny\color{gray},
    stepnumber=1,
    numbersep=5pt,
    frame=single,
    rulecolor=\color{black!30},
    backgroundcolor=\color{white!97!gray},
    breaklines=true,
    captionpos=b,
    commentstyle=\color{green!50!black},
    keywordstyle=\color{blue},
    stringstyle=\color{orange}
}

% ===============================
% Speech-to-Text (Whisper)
% ===============================
\section{Speech-to-Text (Whisper)}

Die folgende Implementierung zeigt die erste Stufe der Pipeline, in der gesprochene Sprache mithilfe des Whisper-Modells in Text transformiert wird.

\begin{lstlisting}[caption={speech_to_text.py – Audio-Transkription mittels Whisper}]
import whisper
import tempfile

model = whisper.load_model("base")

def transcribe_audio(uploaded_file):
    with tempfile.NamedTemporaryFile(delete=False, suffix=".wav") as tmp:
        tmp.write(uploaded_file.read())
        tmp_path = tmp.name

    result = model.transcribe(tmp_path)
    return result["text"]
\end{lstlisting}

% ===============================
% NLP Analyse (spaCy)
% ===============================
\section{Natural Language Processing}

In diesem Schritt erfolgt die linguistische Analyse des transkribierten Textes mithilfe von spaCy. Der Text wird tokenisiert und grammatikalisch annotiert.

\begin{lstlisting}[caption={nlp_analysis.py – Tokenisierung und Part-of-Speech Tagging}]
import spacy

nlp = spacy.load("de_core_news_sm")

def analyze(text):
    doc = nlp(text)

    tokens = [{"text": t.text, "pos": t.pos_} for t in doc]
    verbs = [t.text for t in doc if t.pos_ == "VERB"]
    nouns = [t.text for t in doc if t.pos_ == "NOUN"]

    return {
        "tokens": tokens,
        "verbs": verbs,
        "nouns": nouns
    }
\end{lstlisting}

% ===============================
% Text-Visualisierung (Verben markieren)
% ===============================
\section{Textvisualisierung}

Zur Verbesserung der Interpretierbarkeit werden erkannte Verben im Text visuell hervorgehoben.

\begin{lstlisting}[caption={nlp_analysis.py – Markierung von Verben im Text}]
def highlight_verbs(text):
    doc = nlp(text)
    out = []

    for token in doc:
        if token.pos_ == "VERB":
            out.append(f"<mark>{token.text}</mark>")
        else:
            out.append(token.text)

    return " ".join(out)
\end{lstlisting}

% ===============================
% Bildgenerierung
% ===============================
\section{Visuelle Ausgabe (Bildsuche)}

Die extrahierten Nomen werden als semantische Suchbegriffe verwendet, um passende Bilder aus einer externen Bildquelle abzurufen.

\begin{lstlisting}[caption={image_generator.py – Generierung semantischer Bildabfragen}]
import requests

def get_image_from_nouns(query):
    url = f"https://source.unsplash.com/800x600/?{query}"
    return url
\end{lstlisting}

% ===============================
% Streamlit Anwendung (Pipeline-Orchestrierung)
% ===============================
\section{Streamlit Anwendung}

Dieses Listing zeigt die vollständige Orchestrierung der KI-Pipeline innerhalb einer interaktiven Webanwendung.

\begin{lstlisting}[caption={app.py – Gesamtpipeline des multimodalen Systems}]
import streamlit as st
from speech_to_text import transcribe_audio
from nlp_analysis import analyze, highlight_verbs
from image_generator import get_image_from_nouns

st.title("Multimodaler Sprach-Assistent")

mode = st.radio("Eingabemodus", ["Text", "Audio"])
text = ""

if mode == "Text":
    text = st.text_input("Eingabe")
else:
    audio = st.file_uploader("Audio hochladen")
    if audio:
        text = transcribe_audio(audio)

if text:
    st.subheader("Linguistische Analyse")
    result = analyze(text)
    st.json(result)

    st.subheader("Verben hervorgehoben")
    st.markdown(highlight_verbs(text), unsafe_allow_html=True)

    if result["nouns"]:
        query = " ".join(result["nouns"][:3])
        st.image(get_image_from_nouns(query))
\end{lstlisting}
\end{appendices}

%% ===============
%%  Verzeichnisse
%% ===============

\singlespacing  

\listofreferences
\listoffigures
\listoftables
\listoflistings
\listofacronyms
\printglossary
\printindex

\normalspacing  


%% =========================================
%%  Selbstständigkeitserklärung & Thesen
%% =========================================

\makedoiw  

\maketheses{
Die Kombination aus Speech-to-Text und Natural Language Processing ermöglicht eine effektive Umwandlung gesprochener Sprache in strukturierte Daten;
Whisper stellt eine robuste Grundlage für die Spracherkennung in KI-Systemen dar;
spaCy eignet sich effizient zur linguistischen Analyse und Wortartenbestimmung;
Die Extraktion von Nomen ermöglicht eine sinnvolle Abbildung auf visuelle Inhalte;
Bildbasierte Rückkopplungssysteme erhöhen das Verständnis sprachlicher Informationen;
Streamlit erlaubt eine schnelle Entwicklung interaktiver KI-Anwendungen;
Die modulare KI-Pipeline verbessert die Erweiterbarkeit multimodaler Systeme erheblich.
}

\end{document}